\documentclass[12pt,a4paper]{article}
\usepackage[margin=2.5cm]{geometry}
\usepackage{amsmath,amssymb}
\usepackage{graphicx}
\usepackage{cite}
\usepackage{hyperref}
\hypersetup{colorlinks=true,linkcolor=blue,citecolor=blue,urlcolor=blue}
\usepackage{setspace}
\onehalfspacing

\title{The Transdimensional Energy Descent Cosmological Program:\\
A Unified Cosmological Framework Based on Dimensional Interface Dynamics}
\author{Yi Dong \\
Jinan Lixin Digital Technology Co., Ltd.}
\date{}

\begin{document}
\maketitle

\begin{abstract}
Mainstream cosmology and particle physics adopt a domain‑separated effective‑theory framework. The $\Lambda$CDM model successfully describes large‑scale cosmic evolution, yet the physical nature of dark matter and dark energy remains unresolved, and the underlying mechanisms for singularities and early inflation are missing. The Standard Model of particle physics unifies three microscopic interactions but cannot incorporate gravity, while the ontological dynamical mechanism for mass‑energy conversion has long been left open. General relativity and quantum mechanics exhibit irreconcilable conflicts at the Planck scale. Leading quantum‑gravity proposals introduce many independent ad‑hoc assumptions and have long lacked exclusive observational predictions that can be directly tested.

This paper presents a unified cosmological research programme built upon high‑dimensional‑energy trans‑dimensional descent and dimensional‑interface tension. Starting from four foundational axioms, this framework defines the speed of light $c$ as the intrinsic interface velocity for high‑dimensional energy descending into three‑dimensional space. Within this programme, mass is not an intrinsic static property of particles. Instead, mass emerges as a steady‑state manifestation after descending energy is retarded by crowding collisions and interactions inside 3‑space. This provides a continuous dynamical interpretation for the mass‑energy relation $E=mc^2$.

In this framework, the Big‑Bang singularity corresponds to a white‑hole channel injecting four‑dimensional energy into three‑dimensional space. In the early universe, energy density within an extremely small volume exceeded the saturation threshold of the three‑dimensional host medium. Dense energy units suffered intensive crowding and collisions; descent intrinsic velocity decayed rapidly, generating bulk matter within an extremely short timescale, corresponding to cosmic inflation. Free descending energy avoiding collisional retardation propagates near the speed of light and forms the cosmic microwave background radiation. Gravitation, electromagnetism and nuclear forces are interpreted as derivative effects of trans‑dimensional interfacial tension emerging across different scales and dimensional combinations. Apparent quantum randomness originates from temporal aliasing caused by the observation‑bandwidth limitation of three‑dimensional measurement. Quantum entanglement corresponds to multiple 3‑dimensional projections of a single high‑dimensional object, naturally evading local‑hidden‑variable and superluminal causal paradoxes.

Furthermore, all microscopic particles can be regarded as extrinsic manifestations of homologous high‑dimensional energy units after undergoing different descent depths, dimensional‑collision differentiation and interface‑threshold filtering. The dimension‑bias effect explains particle‑antiparticle pair production. Non‑linear interface‑threshold imagery accounts for discrete quantum numbers, decay‑channel constraints and the physical origin of the three generations of fermions.

This work delivers multiple testable observational predictions together with a full set of first‑order phenomenological dynamical equations, guiding future quantitative validation and discriminative experiments. This programme is an ontological unified framework and does not negate the computational validity of existing physical theories within their respective domains. It aims to connect micro‑, macro‑ and cosmological scales with fewer independent ad‑hoc postulates.

This document is an open‑discussion draft (v0.2 revised). Some equations are phenomenological placeholders; their parameters and asymptotic behaviours will be refined in subsequent revisions. Revisions and discussion materials are hosted in an open‑source repository:
\url{https://github.com/nell332/transdimensional-cosmology-program.git}. Version records are archived on Zenodo:
\url{https://doi.org/10.5281/zenodo.23765255}.

\textbf{Keywords}: trans‑dimensional dynamics; cosmic origin; mass‑energy conversion; quantum entanglement; dark‑matter effect; dimensional‑interface tension; particle homology; descent depth; interface threshold
\end{abstract}

\section{Introduction}
\subsection{Research background and open problems}
Over the past half‑century, cosmology and particle physics have achieved landmark progress. The $\Lambda$CDM cosmological model precisely fits large‑scale cosmic‑structure evolution through high‑precision CMB observations by the Planck satellite \cite{planck2018}. The Standard Model of particle physics unifies electromagnetic, strong and weak interactions; all its predicted particles have obtained experimental confirmation \cite{oerter2006}.

Nevertheless, the existing theoretical system still possesses deep structural limitations:
\begin{enumerate}
\item \textbf{Theoretical fragmentation}: General relativity describes macroscopic gravitation while quantum mechanics governs microscopic particles. They produce fundamental conflicts at the Planck scale and black‑hole singularities. No consensus quantum‑gravity scheme has been established \cite{weinberg1972,rovelli2004}.
\item \textbf{Abundant independent parameters}: To match observational data, the $\Lambda$CDM model introduces dark matter, dark energy and cosmic inflation. Their microscopic origins and underlying dynamical mechanisms lack first‑principle explanations \cite{peebles1993}.
\item \textbf{Unresolved underlying mechanisms}: Special relativity gives the mass‑energy relation $E=mc^2$, yet it does not explain the dynamical process of mutual conversion between mass and energy. The quantum‑measurement problem and non‑locality of entanglement sustain multiple competing interpretations without a unified ontological picture \cite{bell1964,zeilinger2010}.
\end{enumerate}

Leading quantum‑gravity approaches exemplified by string theory prioritise mathematical self‑consistency and symmetry. They introduce extra dimensions, supersymmetry and compactification, but consistently fail to deliver exclusive observationally‑testable predictions. Their physical reality remains widely debated \cite{smolin2006,witten1996}.

\subsection{Research strategy of this work}
Instead of adding corrective terms atop the Standard Model, this paper constructs an axiom‑minimal underlying dynamical framework. It takes high‑dimensional‑energy trans‑dimensional migration, descent and retardation as core mechanisms, incorporating gravitation, microscopic interactions, mass‑energy conversion, quantum effects, cosmic origin and dark‑matter effects within one ontological picture.

This work is a unifying research programme rather than a quantitatively verified mature physical theory. Its core merit lies in resolving multiple logical fractures in existing theoretical systems with fewer independent assumptions and providing a unified underlying picture for subsequent quantitative modelling and experimental testing.

> Declaration: The term ``dimension'' within this programme is first defined as an \textbf{energy‑coupling interface dimension}. It is not identical to compact spatial coordinate dimensions in Kaluza–Klein theory. Whether it can be further interpreted as geometric spatial dimensions is left for future trans‑dimensional metric modelling.

\section{Foundational Axioms and Core Phenomenological Dynamics}
All derivations in this programme rest upon four foundational axioms, with no further ad‑hoc postulates added.

\subsection{Foundational axioms}
\noindent\textbf{Axiom 1. Time‑forward directionality and entropy‑increase principle}\\
Time is intrinsically irreversible. Entropy increase serves as the underlying constraint for energy‑matter transformation. Evolution modes violating time‑forward directionality and entropy increase cannot stably persist within our universe \cite{landau1980}.

\noindent\textbf{Axiom 2. Layered dimensions and trans‑dimensional migration}\\
The universe possesses a layered‑dimensional structure. Higher‑dimensional domains are not subject to work constraints from three‑dimensional space. Energy can undergo trans‑dimensional migration and descent. Descent from higher into lower dimensions possesses an intrinsic interface velocity. Generation and maintenance of dimensional channels must satisfy corresponding energy‑coupling conditions.

\noindent\textbf{Axiom 3. Three‑dimensional space energy‑saturation threshold}\\
Unit volume of three‑dimensional space has an upper saturation energy‑density limit $\rho_{\text{sat}}$ for accommodating descending high‑dimensional energy. When local energy density exceeds this threshold, energy units undergo crowding, collisions and mutual retardation. Each individual energy unit carries fixed intrinsic energy $\varepsilon_0$. Number density and energy density satisfy $\rho = n\varepsilon_0$, where $n$ denotes the number density of energy units.

\noindent\textbf{Axiom 4. Global energy conservation}\\
Total cosmic energy is conserved. Energy cannot be created nor annihilated out of nothing; it merely transforms between different forms and different dimensional‑coupling states \cite{noether1918}.

\subsection{Core physical‑quantity definitions and descent‑dynamics equations}
Define the \textbf{descent intrinsic velocity} $v_d$: an interface intrinsic parameter for high‑dimensional native energy descending into three‑dimensional space.

\textbf{Important distinction}: $v_d$ is an interfacial‑dynamics parameter internal to this programme. It is \textit{not} equal to the coordinate motion velocity of particles inside three‑dimensional spacetime. Their mapping transformation remains to be completed in future work.

When energy suffers no three‑dimensional retardation, the descent intrinsic velocity reaches its maximum value $c$, the vacuum speed of light.

Total energy of a single energy unit decomposes into descent‑kinetic component and rest‑energy component, obeying global conservation:
\begin{equation}
E = E_k + E_m = \text{constant}
\end{equation}
The descent‑kinetic component adopts a first‑order phenomenological form within this programme: $E_k=\frac12 m_0 v_d^2$. This describes energy partitioning inside the trans‑dimensional interface and is not directly equivalent to three‑dimensional relativistic kinetic energy. The rest‑energy component reads $E_m = m(v_d) c^2$. Combining energy conservation yields the phenomenological correspondence between rest mass and descent velocity:
\begin{equation}
m(v_d)=m_0\left(1-\frac{v_d^2}{2c^2}\right)
\end{equation}

Equation (2) is a first‑order phenomenological expansion. For $v_d\ll c$, it reduces to steady‑state rest mass $m_0$ describing ordinary matter. For $v_d\to c$, the first‑order form gives $m\to m_0/2$. To enforce vanishing rest mass at $v_d=c$ for photon‑like pure‑radiation states, we adopt the second‑order corrected form satisfying two‑endpoint boundary conditions:
\begin{equation*}
m(v_d)=m_0\sqrt{1-\frac{v_d^2}{c^2}}
\tag{2b}
\end{equation*}
Equation (2b) expands to $m_0\left(1-v_d^2/(2c^2)\right)$ for $v_d\ll c$, consistent with the low‑order limit of Eq.(2), and strictly yields $m=0$ at $v_d=c$. All subsequent qualitative conclusions requiring vanishing photon rest‑mass use Eq.(2b). Equation (2) is retained as a low‑order approximation to illustrate the monotonic trend of mass varying with descent velocity.

Evolution of descent velocity under three‑dimensional collisional retardation satisfies the first‑order dynamical equation:
\begin{equation}
\frac{dv_d}{dt}=-\alpha\, n\, v_d
\end{equation}
where $\alpha$ is the three‑dimensional collisional‑retardation coefficient (phenomenological fit parameter), and $n$ is local number density of energy units.

Combining Eq.(2b) and Eq.(3), exponential decay of $v_d$ from its initial value toward zero corresponds to rest‑mass growing from zero up to $m_0$, covering the full mass‑generation interval. Using only the first‑order Eq.(2) compresses the mass range between $m_0/2$ and $m_0$, motivating the second‑order correction. The inertia coefficient $m_0$ in $E_k$ is a low‑order phenomenological choice. In higher‑order revisions, inertia‑coefficient dependence on $v_d$ can be absorbed into corrective terms without altering the core monotonic picture of mass versus descent velocity.

\section{Cosmic Origin and Early Evolution: White‑Hole Singularity and Inflation Mechanism}
\subsection{Ontology of the Big Bang: white‑hole energy outpouring}
Within this programme, the Big‑Bang singularity is essentially a white‑hole channel through which four‑dimensional native energy injects into three‑dimensional space. In the very early universe, the white‑hole poured high‑flux descending high‑dimensional energy into a tiny volume of three‑dimensional space. Initial energy density greatly exceeded the three‑dimensional saturation threshold $\rho_{\text{sat}}$.

\subsection{Bulk matter generation and cosmic inflation}
Under Axiom 3, when energy density surpasses the saturation threshold, densely‑packed energy units undergo frequent collisions. Descent intrinsic velocity decays drastically within an extremely short time interval, corresponding to bulk rest‑mass generation. The phenomenological equation for particle number‑density production rate reads:
\begin{equation}
\frac{dn_m}{dt}=k\cdot n^2\cdot\left(1-\frac{n_{\text{sat}}}{n}\right)\cdot H(t)^{-1}
\end{equation}
where $k$ denotes collision‑solidification efficiency coefficient, and $H(t)$ is the Hubble parameter describing dilution effects of spatial expansion \cite{friedmann1922,hubble1929}.

The factor $H(t)^{-1}$ in Eq.(4) captures production enhancement in the high‑density early phase before spatial expansion sufficiently dilutes energy. It applies for $n\gg n_{\text{sat}}$ and non‑vanishing $H(t)$. In the static or extreme deceleration limit $H(t)\to0$, the phenomenological form needs regularisation:
\begin{equation*}
\frac{dn_m}{dt}=k\cdot n^2\cdot\left(1-\frac{n_{\text{sat}}}{n}\right)\cdot\frac{H_0}{H(t)+H_0}
\tag{4b}
\end{equation*}
$H_0$ is the present‑day Hubble constant. For $H\gg H_0$, Eq.(4b) recovers the $H^{-1}$ behaviour of Eq.(4) (the constant factor $H_0$ can be absorbed into $k$). At $H\to0$, it delivers finite production rates and avoids divergence. Both forms yield vanishing production rate as $n\to n_{\text{sat}}$, consistent with the observation that large‑scale matter creation no longer occurs in the present‑day universe.

For $n\gg n_{\text{sat}}$, production rate scales with density‑squared, describing inflation‑like matter creation in the high‑density early cosmos. As $n\to n_{\text{sat}}$, production rate approaches zero and the universe enters gravity‑dominated structural evolution.

This mechanism directly explains why most cosmic matter formed in the very early epoch, while modern space no longer generates bulk new matter: spatial expansion dilutes energy density and removes conditions for crowding‑collision solidification.

Regarding the horizon problem: before bulk matter generation, different three‑dimensional spatial regions maintain trans‑dimensional connectivity through unified high‑dimensional‑object temporal evolution. High‑flux energy injected by the white hole lets all domains of the observable universe share common high‑dimensional evolutionary phases within an extremely short interval, supplying a causal foundation for large‑scale CMB homogeneity. Quantitative causal‑connectivity analysis requires trans‑dimensional‑geometry metric equations and is deferred to future work \cite{guth1981,linde1982}.

\subsection{Origin of the cosmic microwave background radiation}
Not all energy pouring out of the white hole participates in crowding collisions and mass generation. Some energy units experience no effective collisional coupling after ejection. Their descent intrinsic velocity remains close to $c$ without decay or solidification. These freely propagating descending energy units constitute the observed cosmic‑microwave‑background radiation (CMB) \cite{penzias1965,smoot1992}.

This mechanism answers the classic cosmological puzzle: why Big‑Bang remnant radiation still propagates near‑light‑speed across billions of years without converting into matter. It corresponds to pure descending states unhindered by retardation and not satisfying mass‑generating collisional‑coupling conditions.

Concerning CMB black‑body spectral properties: as freely descending energy travels through three‑dimensional space, cosmic expansion stretches its wavelength and dilutes energy density, producing continuous redshift spectra. In the pre‑recombination epoch, weak trans‑dimensional coupling exists between freely descending energy and solidified matter, sufficient to bring them into statistical equilibrium and endow free descending energy with black‑body spectral shape. After recombination this coupling breaks, and free‑descending energy propagates outward with approximately black‑body spectra until today. In addition, statistical distributions of high‑dimensional native energy upon injection intrinsically possess near‑thermal spectral features, consistent with high‑flux white‑hole injection. Rigorous spectral derivation is deferred to future quantitative work \cite{planck2018,komatsu2011}.

Note that the assumption of ``weak trans‑dimensional coupling between free descending energy and solidified matter'' is not among the four foundational axioms and constitutes an open problem for this draft. Future revisions will attempt to derive coupling strength and form from the dimensional‑channel‑coupling condition within Axiom 2.

\section{Unified Explanation of Fundamental Interactions: Derived Trans‑Dimensional‑Tension Effects}
Within this programme, the four fundamental interactions are not mutually independent primitive forces. Instead, they are derivative manifestations of trans‑dimensional interfacial tension under different dimensional combinations and spatial scales.

Strength differences among four interactions correspond to differences in channel width and interfacial tightness for distinct dimensional‑interface combinations. Gravitation corresponds to 3D$\to$2D channels with extremely narrow cross‑section and low leakage rate. Only large‑scale mass aggregation can produce observable deformation, making gravity the weakest interaction. The nuclear force corresponds to tight 4D‑3D interface contact at ultra‑short distances; tension acts directly without dilution and yields the strongest micro‑scale force. Electromagnetism is 4D$\to$3D trans‑dimensional transition tension with intermediate strength, exhibiting attraction and repulsion polarities associated with charge signs. This magnitude narrative supplies a unified channel‑image ordering for interaction strengths. Quantitative strength‑ratio computation is left for future work \cite{griffiths2005,perkins2000}.

\subsection{Gravitation: 3‑to‑2 dimensional trans‑dimensional descent tension}
Gravitation represents spatial distortion caused by three‑dimensional energy descending toward two‑dimensional domains. Define spatial deformation variable $\varepsilon$. It satisfies a Poisson‑type phenomenological equation with local rest‑mass density $\rho_m$:
\begin{equation}
\nabla^2 \varepsilon = 4\pi G_\perp \rho_m
\end{equation}
$G_\perp$ denotes the trans‑dimensional‑tension coefficient. Deformation $\varepsilon$ maps linearly to gravitational potential $\Phi$, i.e. $\Phi\propto\varepsilon$. Under weak‑field, low‑velocity, small‑deformation approximations this reduces to the Newtonian Poisson equation and matches everyday observations. Higher total solidified matter density per unit space amplifies trans‑dimensional descent distortion and enhances gravitational effects \cite{chandrasekhar1965}.

\subsection{Electromagnetism: four‑to‑three dimensional trans‑dimensional transition tension}
Electrons are incompletely frozen high‑dimensional‑energy projections. They intrinsically tend to transition backwards toward four‑dimensional domains and generate reverse trans‑dimensional gravitation, producing attraction between opposite charges. Meanwhile, saturation‑threshold constraints of three‑dimensional space yield saturation repulsion when similar high‑dimensional projections accumulate, corresponding to like‑charge repulsion \cite{jackson1998}.

\subsection{Nuclear force: short‑range strong‑coupling tension of the 4D‑3D interface}
The nuclear force is strong‑binding tension from tight 4D‑3D interfacial coupling at microscopic distances. Its radial dependence adopts a Yukawa‑like potential:
\begin{equation}
V(r)=-V_0 \cdot \frac{e^{-r/r_0}}{r/r_0}
\end{equation}
where $r_0$ is the trans‑dimensional critical coupling distance (on the order of $10^{-15}\,\text{m}$), and $V_0$ is potential‑well depth. When inter‑particle separation falls below this critical scale, the trans‑dimensional‑projection interface fits tightly, coupling conditions are satisfied and strong‑binding interactions emerge. Once separation exceeds $r_0$, trans‑dimensional coupling decays exponentially and interactions vanish rapidly. This reproduces experimental signatures of the nuclear force: short range, strong attraction and saturation behaviour \cite{yukawa1935,halzen1984}.

\subsection{Black holes, white holes and dark‑matter effects}
\begin{itemize}
\item \textbf{Black hole}: stable channel for three‑dimensional energy descending into two‑dimensional domains. Sustained material inflow is required to stabilise the interface. If material supply ceases, the dimensional channel gradually dissipates \cite{hawking1974,chandrasekhar1983}.
\item \textbf{White hole}: injection channel for four‑dimensional energy entering three‑dimensional space. The Big‑Bang singularity constitutes a cosmic‑scale white‑hole event. Black holes and white holes belong to distinct channel categories; there exists no topological ``black‑hole‑to‑white‑hole'' direct passage.
\item \textbf{Dark‑matter effect}: originates from residual trans‑dimensional 3D‑2D distortion tension left over by extinguished ancient black holes. Residual gravitational potential obeys the decay relation:
\begin{equation}
\Phi(r,t)=\Phi_0 \cdot \frac{e^{-t/\tau}}{r}\cdot e^{-r/r_d}
\end{equation}
where $\tau$ is trans‑dimensional‑tension relaxation time, and $r_d$ is spatial‑decay characteristic scale for tension. Dark‑matter signals concentrate in regions that once hosted abundant black‑hole formation. Cosmic voids possess weaker residual tension and lower dark‑matter abundance \cite{rubin1980,bertone2005}.
\end{itemize}

Equation (7) describes gravitational potential left by a single ancient‑black‑hole remnant source inside three‑dimensional space. Total galactic‑scale potential is the integral superposition from many independent remnant sources. Under approximations of uniform remnant spatial distribution and observation scale $r\ll r_d^{\text{gal}}$, total potential takes the logarithmic approximate form:
\begin{equation*}
\Phi_{\text{gal}}(r)\approx\Phi_c \ln\left(\frac{r}{r_c}\right)
\tag{7b}
\end{equation*}
$\Phi_c$ and $r_c$ are constants determined by remnant number density and single‑source intensity. Logarithmic potential yields rotation‑speed relation $v^2 = r\,d\Phi_{\text{gal}}/dr \approx \Phi_c$, producing approximately flat rotation curves inside galactic scales. The exponential cutoff factor $e^{-r/r_d}$ in Eq.(7) describes tension decay at galactic peripheries and eventual drop‑off of rotation curves. In addition, $\tau$ in factor $e^{-t/\tau}$ must be much larger than cosmic age, so that remnants barely decay over cosmic history, consistent with high‑redshift dark‑matter observations. Weak time‑evolution differences among remnants of different ages provide further test directions.

Concerning dimensional‑ladder boundaries: this programme postulates three‑dimensional space as the lowest‑level layer capable of stably hosting accumulated matter. Two‑dimensional and lower‑dimensional domains only exist in forms of transient deformation, remnant tension or channel cross‑sections and cannot stably accumulate matter. Sustained energy descent toward two‑dimensional domains demands extreme‑density material inflow (e.g.~black holes). Channels dissolve once material supply terminates. Therefore the dimensional ladder ``bottoms out'' at three dimensions; stable ``2‑dimensional matter'' or ``1‑dimensional matter'' forms do not exist.

The weak interaction is only qualitatively conjectured as metastable‑couple dissociation processes on trans‑dimensional interfaces. Corresponding phenomenological potentials, decay channels and parity‑violation signatures constitute primary future‑modelling targets. No quantitative forms are supplied within this paper \cite{weinberg1967,wilczek2004}.

\section{Trans‑Dimensional‑Projection Interpretation of Quantum Phenomena}
\subsection{Apparent randomness: temporal aliasing from limited observation bandwidth}
The full ontological being of quantum objects resides in high‑dimensional space. Its state evolution is deterministic and continuous in time and possesses no intrinsic probabilistic character.

All three‑dimensional observational means are constrained by three‑dimensional causality and signal velocity, giving an intrinsic minimum time slice, i.e.~an upper bound $f_{\text{obs}}$ for observational sampling frequency.

\textbf{Remark}: The temporal‑aliasing analogy used here is an ontological model inside this framework. It is not directly copied from classical signal‑processing ADC sampling. The three‑dimensional observation‑bandwidth upper bound originates from light‑speed causality constraints rather than instrument hardware sampling rates.

Intrinsic high‑dimensional evolution frequency $f_d$ is far higher than three‑dimensional observational sampling frequency, satisfying the Nyquist undersampling condition:
\begin{equation}
f_{\text{obs}} < 2 f_d
\end{equation}

Temporal aliasing then occurs. Misalignment between sampling phases and high‑dimensional‑evolution phases produces apparent randomness in measurement outcomes. The quantum uncertainty principle is essentially an aliasing effect from insufficient observation bandwidth, not intrinsic indeterminacy of objects \cite{shannon1949,zeilinger2000}.

Let $\omega_d$ denote angular frequency of high‑dimensional ontological evolution and $\omega_s$ the observational‑sampling angular frequency. Apparent aliased frequency after observation reads:
\begin{equation}
\omega_{\text{app}}=\left|\omega_d - k\,\omega_s\right|
\end{equation}
where $k$ is integer satisfying aliasing conditions. Apparent probability distribution of observational outcomes is determined by phase statistics averaged over discrete sampling points along deterministic high‑dimensional trajectories:
\begin{equation}
P(x)=\frac{1}{T_s}\int_{0}^{T_s}\delta\big(x-x_d(t)\big)\,dt
\end{equation}
$T_s$ stands for sampling period; $x_d(t)$ is deterministic high‑dimensional ontological trajectory.

Equation (10) gives apparent probability distribution for single‑projection‑channel cases. For multi‑spatially‑separated‑channel experiments such as double‑slit interference, observational probabilities must incorporate phase factors and cross‑terms from high‑dimensional trajectories of each channel:
\begin{equation*}
P_{\text{multi}}(x)\propto\left|\sum_j A_j e^{i\phi_j(x)}\right|^2
\tag{10b}
\end{equation*}
$A_j$ is amplitude for the $j$-th projection channel; $\phi_j(x)$ accumulates high‑dimensional‑trajectory phase at observation position $x$. Cross‑term $2A_1A_2\cos(\phi_1-\phi_2)$ generates interference fringes. For single‑channel scenarios, Eq.(10b) reduces to time‑averaged classical form Eq.(10). For multi‑channel scenarios, phase correlations of high‑dimensional trajectories are projected into three‑dimensional interference patterns via sampling processes. Quantum interference is therefore not intrinsic wave‑like behaviour of objects but apparent manifestation of multi‑channel high‑dimensional‑trajectory phase correlations under finite‑bandwidth sampling \cite{feynman1963,wheeler1983}.

It should be honestly noted that the phase‑superposition form in Eq.(10b) borrows probability‑amplitude superposition rules from quantum mechanics. This framework has not yet strictly derived the $A_j e^{i\phi_j(x)}$ structure starting from high‑dimensional deterministic‑trajectory sampling. This represents the core unresolved gap within the quantum section. Equation (10b) serves here as a phenomenological placeholder: if multi‑projection‑channel phase correlations exist, interference patterns shall take this form. First‑principle derivation is deferred to future work.

\subsection{Quantum‑entanglement mechanism: multiple projections from a single high‑dimensional object}
A core postulate of this programme: there do not exist two mutually independent entangled particles. Observed quantum‑entanglement phenomena correspond to multiple spatially‑separated images projected by one single high‑dimensional object onto three‑dimensional spacetime. Two projections far apart within three‑dimensional view share identical high‑dimensional‑object evolutionary timelines and possess intrinsic correlation. No superluminal signal needs to propagate inside three‑dimensional space; three‑dimensional light‑speed causality constraints remain unbroken.

This mechanism naturally evades local‑hidden‑variable restrictions of Bell’s theorem. Derivation of Bell inequalities presupposes ``two subsystems are locally independent entities''. Since the two projections are not independent entities within our framework, that premise does not apply. No conflict arises with Bell‑experiment violations of Bell inequalities \cite{bell1964,aspect1982}.

It should be noted that in discussions allowing superdeterminism, Bell’s theorem still requires measurement‑choice independence from hidden variables (free‑choice condition). Within this framework, measurement bases are selected by three‑dimensional observers and are operationally independent from high‑dimensional‑object evolutionary timelines, satisfying this condition. No three‑dimensional‑space signal passes between projections; correlation origins are placed inside unified high‑dimensional‑object evolution. More rigorous ontological argumentation for the free‑choice condition awaits further work. This paper adopts the operational standpoint that observer basis‑selection behaves as independent from high‑dimensional‑object evolution at three‑dimensional descriptive level.

\subsection{Testing proposal}
Precisely record timestamps, measurement bases and outcomes for every run in entanglement experiments. Perform spectral‑search analysis on time‑series datasets, searching for reproducible intrinsic high‑dimensional‑evolution frequencies that can predict subsequent measurement results. Discovery of stable intrinsic frequencies would supply direct supporting evidence for this model.

\section{Homologous Origin of Microscopic Particles: Descent‑Depth and Dimensional‑Interface Threshold Effects}
Within the picture of this programme, diverse microscopic particles are not a set of mutually‑independent fundamental building blocks pre‑existing in the universe. Instead they are different steady‑state appearances produced by one single species of primordial high‑dimensional‑energy unit after trans‑dimensional descent, three‑dimensional collisions and non‑linear filtering at dimensional interfaces. Mass is directly equivalent to the \textbf{descent depth} of energy units: descent depth equals zero corresponds to energy fully unretarded by three‑dimensional media and zero rest‑mass. Deeper descent means stronger three‑dimensional collisional retardation and larger resulting rest‑mass.

\subsection{Dimensional‑bias origin of particle‑antiparticle pairs}
During three‑dimensional crowding collisions, high‑dimensional‑energy units undergo dimensional‑attribute differentiation, analogous to magnets which always keep two poles no matter how finely you cut them. After collisions, some energy units retain more four‑dimensional‑side attributes, while others develop stronger three‑dimensional‑side bias. This produces paired particles and corresponding antiparticles. Taking electron‑positron as an example: the negative electron preserves more four‑dimensional‑attribute bias, whereas the positron manifests stronger three‑dimensional‑attribute bias.

When particles and antiparticles carrying opposite dimensional biases encounter each other, their dimensional tendencies cancel mutually. Units revert back to un‑solidified primitive descending energy, i.e.~particle‑antiparticle annihilation occurs. This picture supplies an ontological explanation for pair production and annihilation based upon dimensional‑interface collisions.

\subsection{Non‑linear interface thresholds: discrete quantum numbers and decay‑channel constraints}
Energy descent into three‑dimensional space is not a smoothly continuously‑tunable process. Dimensional interfaces possess intrinsic non‑linear state thresholds analogous to phase transitions in matter. Intermediate descent states between distinct steady‑states can in principle exist, yet sustaining such intermediate configurations demands extremely precise local conditions which natural three‑dimensional cosmic environments can barely maintain stably. Once conditions deviate from threshold narrow gaps, systems jump directly toward a small set of discrete steady‑states permitted by interfaces.

Therefore experimentally observed quantum numbers such as charge and spin take strictly discrete values. Intermediate transition configurations cannot survive stably during particle decay, and decay proceeds through a restricted set of allowed channels. This does not mean intermediate states are impossible in principle, but rather that three‑dimensional cosmic environments cannot preserve intermediate configurations for long times. Resonance‑threshold mechanisms at dimensional interfaces filter out only a small subset of particle steady‑states capable of long‑term survival.

\subsection{Three generations of fermions: different descent‑depth scales inside one threshold interval}
Electron, muon and tau (three fermion generations) are homologous high‑dimensional‑energy units lying inside one broad permitted‑threshold interval of dimensional interfaces. Constrained by interface resonance, their core quantum features such as spin and charge remain identical. Nevertheless they differ in \textbf{descent‑depth scales}: tau possesses the deepest descent depth, highest mass and poorest stability with fastest decay rates; muon sits intermediate; the electron has the shallowest descent depth and can survive stably over cosmic timescales.

Three fermion generations are not three completely independent fundamental‑particle species. They correspond to distinct steady‑states of homologous energy at different descent depths within one threshold band. The two heavier generations merely cannot persist long inside natural cosmic environments.

This section belongs to qualitative phenomenological extension. Mathematical forms for thresholds, quantitative functional relations connecting descent‑depth with quantum numbers and mass, together with why exactly three fermion generations appear, await dynamical‑equation construction in future revisions. Meanwhile this imagery suggests that transient short‑lived intermediate resonance states might in principle be excited and observed under specific high‑energy‑collision windows, constituting one potential discriminative direction for future investigation.

\section{Observational Predictions and Discriminative Test Directions}
If this programme merely explains already‑known phenomena, it remains only an interpretive framework. To become testable physical theory it must deliver exclusive predictions different from the Standard Model. Three discriminative directions are proposed below.

\subsection{Prediction 1: periodic modulation in high‑time‑resolution quantum measurements}
Quantum randomness originates from observational temporal aliasing. As observational time resolution continuously improves and approaches intrinsic three‑dimensional minimal time‑slice limits, measurable periodic‑modulation signatures will appear within quantum statistical distributions. Standard quantum mechanics holds probability rules independent of measurement time‑resolution and predicts no such effect. Observation of this modulation would provide direct experimental evidence supporting this programme.

\subsection{Prediction 2: correlation between dark‑matter distribution and ancient‑black‑hole relics}
Dark‑matter effects originate from residual trans‑dimensional tension left by ancient extinguished black holes. Therefore dark‑matter‑density peaks shall positively correlate with historical regions of intensive black‑hole formation. Cosmic large‑scale void regions will exhibit significantly lower dark‑matter abundance compared to uniform‑distribution expectations of $\Lambda$CDM. This prediction can be examined via next‑generation weak‑lensing and galaxy‑survey datasets \cite{kaiser1992,des2021}.

\subsection{Prediction 3: systematic deviations within primordial‑nucleosynthesis element‑abundance fits}
Based upon the crowding‑retardation matter‑creation model, quantitative fits for primordial‑nucleosynthesis element abundances will produce systematically distinguishable deviations from standard $\Lambda$CDM results. Discrimination can be performed using high‑precision primordial‑element‑abundance observations \cite{burbidge1957,coc2015}.

\section{Discussion and Outlook}
\subsection{Strengths of the framework}
\begin{enumerate}
\item \textbf{Axiom‑minimal and unified narrative}: Starting with only four foundational axioms, it connects cosmic origin, mass‑energy conversion, fundamental interactions, quantum effects and dark‑matter effects. Qualitative imagery further extends to particle homology, particle‑antiparticle origin and three‑generation‑fermion phenomenology, without introducing abundant independent ad‑hoc hypotheses.
\item \textbf{Resolving classical paradoxes}: It supplies underlying dynamical interpretations for mass‑energy conversion. Projection‑based imagery resolves entanglement non‑locality and superluminal‑causal paradoxes instead of setting aside contradictions via domain‑of‑validity demarcation.
\item \textbf{Compatibility with existing physics}: It does not negate computational accuracy of relativity, quantum mechanics and $\Lambda$CDM inside their respective domains. They are treated as effective approximations of this framework under specific conditions.
\end{enumerate}

\subsection{Limitations and follow‑up work}
This research remains at the ontological‑programme stage. Complete sets of quantitative dynamical equations are not yet established, and substantial work separates it from mature physical theory:
\begin{enumerate}
\item Construct unified trans‑dimensional‑geometry temporal differential equations for unified description of four fundamental interactions; complete quantitative relations for threshold conditions, descent‑depth and physical quantities inside the particle section; explain the origin of exactly three fermion generations.
\item Strictly distinguish descent intrinsic velocity and three‑dimensional coordinate motion velocity, build transformation rules consistent with experimental consequences of Lorentz transformations.
\item Improve quantitative modelling for early‑universe inflation and element nucleosynthesis, perform fitting comparisons against CMB observational datasets.
\item Design high‑time‑resolution quantum experiments to test periodic‑modulation predictions; assess feasible energy windows for searching transient intermediate descent‑state resonances at high‑energy colliders.
\end{enumerate}

This programme is currently an open‑discussion draft (v0.2 revised). Some equations are phenomenological placeholders. Their parameters and asymptotic behaviours will be refined in subsequent revisions. Updated revisions and discussion materials are maintained within the open‑source repository. Versioned DOI and repository links are preserved within Zenodo records.

\subsection{Philosophical boundary}
Even if quantitative modelling is completed and partial observational tests are satisfied in the future, experimental data alone cannot ``prove'' the high‑dimensional‑descent picture as objective cosmic truth. Within philosophy‑of‑science theory‑equivalence problems, multiple ontologically‑distinct models can yield identical observational predictions. The core function of experiments is falsification rather than uniquely confirming one particular ontological picture.

The principal merit of this programme lies in offering an alternative ontological picture with fewer independent ad‑hoc assumptions and high internal narrative consistency, supplying an alternative research path for fundamental physics distinct from mainstream paradigms \cite{popper1959,kuhn1962}.

\begin{thebibliography}{99}
\bibitem{planck2018} Planck Collaboration. Planck 2018 results. VI. Cosmological parameters[J]. Astronomy \& Astrophysics, 2020, 641: A6.
\bibitem{oerter2006} Oerter R. The Theory of Almost Everything: The Theory of Almost Everything: The Unsung Triumph of Modern Physics[M]. Penguin, 2006.
\bibitem{weinberg1972} Weinberg S. Gravitation and Cosmology: Principles and Applications of the General Theory of Relativity[M]. Wiley, New York, 1972.
\bibitem{rovelli2004} Rovelli C. Quantum Gravity[M]. Cambridge University Press, Cambridge, 2004.
\bibitem{peebles1993} Peebles P J E. Principles of Physical Cosmology[M]. Princeton University Press, Princeton, 1993.
\bibitem{bell1964} Bell J S. On the Einstein Podolsky Rosen paradox[J]. Physics Physique Fizika, 1964, 1(3):195‑200.
\bibitem{zeilinger2010} Zeilinger A. Dance of the Photons: From Einstein to Quantum Teleportation[M]. Farrar, Straus and Giroux, New York, 2010.
\bibitem{smolin2006} Smolin L. The Trouble with Physics: The Rise of String Theory, the Fall of a Science[M]. Houghton Mifflin Harcourt, Boston, 2006.
\bibitem{witten1996} Witten E. String theory dynamics in various dimensions[J]. Nuclear Physics B, 1995, 443(1‑2):85‑126.
\bibitem{landau1980} Landau L D, Lifshitz E M. Statistical Physics Part 1[M]. Pergamon Press, Oxford, 1980.
\bibitem{noether1918} Noether E. Invariante Variationsprobleme[J]. Nachrichten von der Gesellschaft der Wissenschaften zu Göttingen, Mathematisch‑Physikalische Klasse, 1918: 235‑257.
\bibitem{friedmann1922} Friedmann A. Über die Krümmung des Raumes[J]. Zeitschrift für Physik, 1922, 10(1):377‑386.
\bibitem{hubble1929} Hubble E P. A relation between distance and radial velocity among extra‑galactic nebulae[J]. Proceedings of the National Academy of Sciences, 1929, 15(3):168‑173.
\bibitem{guth1981} Guth A H. Inflationary universe: A possible solution to the horizon and flatness problems[J]. Physical Review D, 1981, 23(2): 347‑356.
\bibitem{linde1982} Linde A D. A new inflationary universe scenario: A possible solution of the horizon, flatness, homogeneity, isotropy and primordial monopole problems[J]. Physics Letters B, 1982, 108(6):389‑393.
\bibitem{penzias1965} Penzias A A, Wilson R W. A measurement of excess antenna temperature at 4080 Mc/s[J]. The Astrophysical Journal, 1965, 142:419.
\bibitem{smoot1992} Smoot G F, et al. Structure in the COBE differential microwave radiometer first‑year maps[J]. The Astrophysical Journal Letters, 1992, 396(1): L1‑L5.
\bibitem{komatsu2011} Komatsu E, et al. Seven‑year Wilkinson Microwave Anisotropy Probe (WMAP) observations: cosmological interpretation[J]. The Astrophysical Journal Supplement Series, 2011, 192(1):18.
\bibitem{griffiths2005} Griffiths D J. Introduction to Elementary Particles[M]. Wiley‑VCH, Weinheim, 2005.
\bibitem{perkins2000} Perkins D H. Introduction to High Energy Physics[M]. Cambridge University Press, Cambridge, 2000.
\bibitem{chandrasekhar1965} Chandrasekhar S. The Mathematical Theory of Black Holes[M]. Oxford University Press, Oxford, 1965.
\bibitem{jackson1998} Jackson J D. Classical Electrodynamics[M]. Wiley, New York, 1998.
\bibitem{yukawa1935} Yukawa H. On the interaction of elementary particles[J]. Proceedings of the Physico‑Mathematical Society of Japan, 1935, 17: 48‑57.
\bibitem{halzen1984} Halzen F, Martin A D. Quarks and Leptons: An Introductory Course in Modern Particle Physics[M]. Wiley, New York, 1984.
\bibitem{hawking1974} Hawking S W. Black hole explosions?[J]. Nature, 1974, 248(5443): 30‑31.
\bibitem{rubin1980} Rubin V C, Ford W K, Thonnard N. Rotation of the Andromeda Nebula from a spectroscopic survey of emission regions[J]. The Astrophysical Journal, 1980, 238:471‑487.
\bibitem{bertone2005} Bertone G, Hooper D, Silk J. Particle dark matter: evidence, candidates and constraints[J]. Physics Reports, 2005, 405(5‑6):279‑390.
\bibitem{weinberg1967} Weinberg S. A model of leptons[J]. Physical Review Letters, 1967, 19(21): 1264‑1266.
\bibitem{wilczek2004} Wilczek F. Origin of mass[J]. Annals of Physics, 2004, 309(2):137‑149.
\bibitem{shannon1949} Shannon C E. Communication in the presence of noise[J]. Proceedings of the IRE, 1949, 37(1): 10‑21.
\bibitem{zeilinger2000} Zeilinger A. The principle of complementarity[J]. Reviews of Modern Physics, 2000, 72(3): 885‑896.
\bibitem{feynman1963} Feynman R P, Leighton R B, Sands M. The Feynman Lectures on Physics[M]. Addison‑Wesley, Reading, 1963.
\bibitem{wheeler1983} Wheeler J A, Zurek W H (Eds.). Quantum Theory and Measurement[M]. Princeton University Press, Princeton, 1983.
\bibitem{aspect1982} Aspect A, Dalibard J, Roger G. Experimental test of Bell’s inequalities using time‑varying analyzers[J]. Physical Review Letters, 1982, 49(25): 1804‑1807.
\bibitem{brassard2000} Brassard G, et al. A review of some aspects of Bell’s theorem[J]. Foundations of Physics, 2000, 30(11): 1601‑1623.
\bibitem{gill2014} Gill R D. Bell’s theorem: what it takes to prove it[J]. Statistical Science, 2014, 29(4): 512‑528.
\bibitem{kaiser1992} Kaiser N. Weak gravitational lensing[J]. Annual Review of Astronomy and Astrophysics, 1992, 30:499‑536.
\bibitem{des2021} Dark Energy Survey Collaboration. Dark Energy Survey Year 3 results: Weak‑lensing and galaxy‑clustering constraints on cosmology[J]. Physical Review D, 2021, 104(12):123529.
\bibitem{burbidge1957} Burbidge E M, Burbidge G R, Fowler W A, Hoyle F. Synthesis of the elements in stars[J]. Reviews of Modern Physics, 1957, 29(4): 547‑650.
\bibitem{coc2015} Coc A, Vangioni‑Flam E. Primordial nucleosynthesis[J]. Annual Review of Nuclear and Particle Science, 2015, 65: 377‑406.
\bibitem{popper1959} Popper K R. The Logic of Scientific Discovery[M]. Routledge, London, 1959.
\bibitem{kuhn1962} Kuhn T S. The Structure of Scientific Revolutions[M]. University of Chicago Press, Chicago, 1962.
\end{thebibliography}

\end{document}
